\appendix

\chaptertitlepage{Dataset Details}
\label{app:datasets}

\section{Dataset Overview}

The dataset has four mutually exclusive classes: Normal, Tuberculosis, COVID-19, and Pneumonia. In the reported run, COVID-19 has 130 images after deduplication by content hash, Normal and Pneumonia are capped at 1{,}200 images each, and training uses inverse-frequency class weights.

\subsection{Data Sources}

\begin{table}[H]
\centering
\caption{Dataset sources, paths, and known public links}
\label{tab:app-datasets}
\footnotesize
\begin{adjustbox}{max width=\textwidth}
\begin{tabular}{@{}P{2.6cm}P{3.4cm}P{3.2cm}P{3.0cm}@{}}
\toprule
\HeaderRow \Head{Dataset} & \Head{Local path} & \Head{Supplies} & \Head{Public link} \\
\midrule
TB Chest X-Ray & \path{TB_Chest_Radiography_Database/} & TB, Normal & Kaggle TB CXR corpus \\
\ZebraRow Covid19-dataset & \path{DenseNet121/Covid19-dataset/} & COVID-19, Normal, Pneumonia & Local COVID folder (317 images on checkout) \\
Chest X-Ray Pneumonia & \path{pneumonia_detection/data/chest_xray/} & Pneumonia, Normal & Kaggle CXR pneumonia corpus \\
\bottomrule
\end{tabular}
\end{adjustbox}
\end{table}

The manifest used for evaluation has 3{,}227 images: 2{,}260 for training, 485 for validation, and 482 for testing, with 130 COVID-19 images after deduplication. Viral Pneumonia images from the COVID source are counted as Pneumonia. Chapter~5 gives the full split and class-by-source tables.

\subsection{Preprocessing and Splits}

Chapter~3 describes the patient-grouped splits (seed 1337), deduplication by content hash, the shared geometric augmentation, and each architecture's input size and preprocessing.

\chaptertitlepage{Model Architecture Summaries}
\label{app:architectures}

\section{Model Training and Implementation}

\begin{itemize}
    \item EfficientNetB3 / DenseNet121: ImageNet base plus shared head (GAP, Dropout, Dense 256, Dropout, Softmax 4); Phase~1 frozen; Phase~2 unfreeze from layers 340 / 300 ($\sim$43\%).
    \item Custom CNN: five-block stack ending in Softmax(4).
    \item Scripts: \path{training/build_manifest.py}, \path{common.py}, \path{models.py}, \path{train.py}, \path{evaluate.py}, \path{leakage_probe.py}; frontend under \path{Frontend/}; weights \path{efficientnetb3_4class.keras}, \path{densenet121_4class.keras}, \path{cnn_4class.keras}.
    \item Run order: build manifest, then train all, then evaluate, optional leakage probe, then \path{Frontend/app.py}. A portable inference bundle under \path{xray-classifier-portable/} already ships the three \texttt{*.keras} weight files for demo without retraining on the local machine.
\end{itemize}

\chaptertitlepage{Alignment with Sustainable Development Goals (SDGs)}
\label{app:sdgs}

\begin{table}[H]
\centering
\caption{Potential educational and research alignment with SDGs (non-clinical)}
\label{tab:app-sdgs}
\footnotesize
\begin{adjustbox}{max width=\textwidth}
\begin{tabular}{@{}cP{3.5cm}P{7.5cm}@{}}
\toprule
\HeaderRow \Head{No.} & \Head{SDG} & \Head{Contribution language} \\
\midrule
1 & SDG 3 Good Health and Well-being & Decision-support \emph{research} for respiratory screening accessibility, with explicit non-diagnostic limits. \\
\ZebraRow 2 & SDG 4 Quality Education & Reproducible teaching artefact for medical AI methodology (fair splits, metrics literacy). \\
3 & SDG 9 Industry, Innovation and Infrastructure & Open software pipeline for multi-model CXR benchmarking. \\
\ZebraRow 4 & SDG 10 Reduced Inequalities & Potential relevance to resource-limited screening contexts as \emph{future clinical research}; not a deployed medical device. \\
\bottomrule
\end{tabular}
\end{adjustbox}
\end{table}

We keep these claims modest, since this FYP is an educational research demonstration built on public data.
